\documentclass[12pt,a4paper]{article}

\usepackage[utf8]{inputenc}
\usepackage[T1]{fontenc}
\usepackage{geometry}
\usepackage{hyperref}
\usepackage{xcolor}
\usepackage{listings}

\geometry{margin=2.5cm}

\lstdefinestyle{code}{
    basicstyle=\ttfamily\small,
    breaklines=true,
    frame=single,
    columns=fullflexible
}

\title{Protocol de Comunicare \\ T17 - Analiza Semantica Cod Sursa}
\author{Proiect PCD}
\date{2026}

\begin{document}

\maketitle
\tableofcontents
\newpage

\section{Introducere}

Acest document descrie protocoalele de comunicare folosite in proiectul T17 - Analiza Semantica Cod Sursa.

Proiectul foloseste trei forme de comunicare:

\begin{itemize}
    \item protocol socket text-based pentru clientii C/Python si admin;
    \item REST API minimal pentru acces HTTP cu raspunsuri JSON;
    \item SOAP API minimal pentru acces HTTP cu raspunsuri XML SOAP.
\end{itemize}

Clientii normali comunica prin:

\begin{lstlisting}[style=code]
127.0.0.1:8080
\end{lstlisting}

Clientul admin comunica prin:

\begin{lstlisting}[style=code]
/tmp/t17_admin.sock
\end{lstlisting}

REST API-ul ruleaza pe:

\begin{lstlisting}[style=code]
http://127.0.0.1:8081
\end{lstlisting}

SOAP API-ul ruleaza pe:

\begin{lstlisting}[style=code]
http://127.0.0.1:8082
\end{lstlisting}

\section{Model de comunicare socket}

Fiecare conexiune socket incepe obligatoriu cu o comanda de autentificare:

\begin{lstlisting}[style=code]
LOGIN <username> <password>
\end{lstlisting}

Dupa autentificare, clientul trimite o comanda principala catre server. Serverul proceseaza comanda si raspunde.

Serverul creeaza cate un proces copil pentru fiecare conexiune acceptata. Socket-urile principale sunt monitorizate cu \texttt{select()}.

\section{Format raspunsuri socket}

\subsection{Raspuns text}

Pentru raspunsuri text, serverul foloseste formatul:

\begin{lstlisting}[style=code]
RESULT <size>
<body>
\end{lstlisting}

\subsection{Raspuns fisier}

Pentru download de raport, serverul foloseste formatul:

\begin{lstlisting}[style=code]
FILE <filename> <size>
<file_content>
\end{lstlisting}

\section{Autentificare}

Comanda de autentificare:

\begin{lstlisting}[style=code]
LOGIN <username> <password>
\end{lstlisting}

Exemple:

\begin{lstlisting}[style=code]
LOGIN user1 pass1
LOGIN admin admin123
\end{lstlisting}

Utilizatorii sunt cititi din:

\begin{lstlisting}[style=code]
config/users.cfg
\end{lstlisting}

Format:

\begin{lstlisting}[style=code]
username:password:role
\end{lstlisting}

\section{Comenzi pentru client normal}

\subsection{UPLOAD}

\begin{lstlisting}[style=code]
UPLOAD <filename> <size>
<file_content>
\end{lstlisting}

Serverul salveaza fisierul in \texttt{uploads/}, ruleaza analyzer-ul si trimite rezultatul inapoi.

\subsection{DOWNLOAD\_REPORT}

\begin{lstlisting}[style=code]
DOWNLOAD_REPORT <report_name>
\end{lstlisting}

Serverul cauta raportul in \texttt{reports/} si il trimite inapoi cu formatul \texttt{FILE}.

\subsection{QUIT}

\begin{lstlisting}[style=code]
QUIT
\end{lstlisting}

\section{Comenzi pentru admin}

Clientul admin foloseste socket UNIX:

\begin{lstlisting}[style=code]
/tmp/t17_admin.sock
\end{lstlisting}

Comenzile admin sunt acceptate doar daca conexiunea vine prin socket UNIX si rolul este \texttt{admin}.

\begin{lstlisting}[style=code]
STATS
LOGS
LIST_UPLOADS
LIST_REPORTS
SERVER_STATUS
CLEAR_LOGS
DELETE_REPORT <report_name>
LIST_USERS
JOBS
QUIT
\end{lstlisting}

\section{REST API}

REST API-ul ruleaza separat:

\begin{lstlisting}[style=code]
python3 rest_api/app.py
\end{lstlisting}

Adresa:

\begin{lstlisting}[style=code]
http://127.0.0.1:8081
\end{lstlisting}

Endpoint-uri disponibile:

\begin{lstlisting}[style=code]
GET /
GET /health
GET /stats
GET /reports
GET /reports/<name>
GET /uploads
GET /jobs
GET /users
GET /logs
GET /ui
\end{lstlisting}

REST API-ul intoarce raspunsuri JSON si citeste date din:

\begin{itemize}
    \item \texttt{logs/stats.txt};
    \item \texttt{logs/jobs.log};
    \item \texttt{logs/server.log};
    \item \texttt{reports/};
    \item \texttt{uploads/};
    \item \texttt{config/users.cfg}.
\end{itemize}

Testare:

\begin{lstlisting}[style=code]
./scripts/test_rest.sh
curl http://127.0.0.1:8081/health
curl http://127.0.0.1:8081/stats
curl http://127.0.0.1:8081/reports
\end{lstlisting}

\section{SOAP API}

SOAP API-ul ruleaza separat:

\begin{lstlisting}[style=code]
python3 soap_api/soap_server.py
\end{lstlisting}

Endpoint SOAP:

\begin{lstlisting}[style=code]
http://127.0.0.1:8082/soap
\end{lstlisting}

Document WSDL:

\begin{lstlisting}[style=code]
http://127.0.0.1:8082/wsdl
\end{lstlisting}

Operatii SOAP suportate:

\begin{lstlisting}[style=code]
Health
Stats
Reports
Uploads
Jobs
Users
Logs
\end{lstlisting}

SOAP API-ul primeste cereri XML in format SOAP Envelope si intoarce raspunsuri XML tot in format SOAP Envelope.

Exemplu cerere SOAP:

\begin{lstlisting}[style=code]
<?xml version="1.0" encoding="UTF-8"?>
<soap:Envelope xmlns:soap="http://schemas.xmlsoap.org/soap/envelope/">
  <soap:Body>
    <Health/>
  </soap:Body>
</soap:Envelope>
\end{lstlisting}

Exemplu raspuns SOAP:

\begin{lstlisting}[style=code]
<?xml version="1.0" encoding="UTF-8"?>
<soap:Envelope xmlns:soap="http://schemas.xmlsoap.org/soap/envelope/">
  <soap:Body>
    <HealthResponse>
      <status>running</status>
      <service>T17 SOAP API</service>
      <host>127.0.0.1</host>
      <port>8082</port>
    </HealthResponse>
  </soap:Body>
</soap:Envelope>
\end{lstlisting}

Testare:

\begin{lstlisting}[style=code]
./scripts/test_soap.sh
curl http://127.0.0.1:8082/wsdl
\end{lstlisting}

\section{Coduri de eroare socket}

Erorile socket sunt trimise ca raspuns \texttt{RESULT}.

Exemple:

\begin{lstlisting}[style=code]
ERROR: login required
ERROR: invalid credentials
ERROR: admin commands require UNIX socket
ERROR: admin permission required
ERROR: unknown command
ERROR: analyzer failed
ERROR: cannot delete report
ERROR: report file not found
\end{lstlisting}

\section{Concluzie}

Protocolul socket este simplu, text-based si usor de testat. Clientii normali folosesc comunicare INET TCP, iar administrarea se face separat prin socket UNIX. In plus, proiectul include REST API minimal cu raspunsuri JSON si SOAP API minimal cu raspunsuri XML/WSDL.

\end{document}